\clearpage\fancyhead[L]{\small JUNO.}\hypertarget{native-335-lower}{}\label{native-335-lower}
\begin{center}\textbf{JUNO.}\par\vspace{4pt}\hrule\vspace{4pt}\fontsize{10}{12}\selectfont Monatliche Correspondenz. Band XXIV. Seite 186\,..\,188. 1811. August.\par\vspace{4pt}\hrule\vspace{4pt}\end{center}

\par\indent Die zum achtenmal vom Herrn Prof. \textsc{Gauss} verbesserten Elemente der \emph{Juno} theilten wir unsern Lesern im Septemberheft 1808 dieser Zeitschrift [S. 307 d. B.] mit, und seitdem konnte eine neue Verbesserung nicht füglich vorgenommen werden, da die fünfte im Januar 1810 eingetretene Opposition der \emph{Juno} nirgends beobachtet wurde. Desto wichtiger war es daher, die diesjährige im April statt findende Opposition nicht zu versäumen; wegen ungemeiner Lichtschwäche des Planeten, war die Beobachtung mit vorzüglichen Schwierigkeiten verknüpft, und wahrscheinlich ist Herr Prof. \textsc{Gauss} der einzige, dem die Beobachtung dieses Gegenscheins gelungen ist. Hier war der Planet in dem so lichtstarken Mittagsfernrohr der Seeberger Sternwarte nie sichtbar, und die wenigen, bei dem äusserst geringen Lichte des Planeten nur am Kreis-Micrometer von Herrn Prof. \textsc{Gauss} angestellten Beobachtungen, sind um so schätzbarer. Zur Vergleichung wurden $\varphi$ Librae und einige andere in der Nähe befindliche Sterne der \emph{Histoire céleste} angewandt, nach deren künftiger schärferer Bestimmung die folgenden Resultate noch einer kleinen Berichtigung bedürftig sein werden.
\Needspace{10\baselineskip}\begin{center}\fontsize{10}{12}\selectfont\setlength{\tabcolsep}{4pt}\renewcommand{\arraystretch}{1.15}
\begin{tabular}{l|c|c|c}
1811 & \shortstack{Mittlere Zeit\\in Göttingen} & \shortstack{Gerade\\Aufsteigung} & \shortstack{Südliche\\Declination} \\
\hline
April 22 & $9^{\mathrm u}\,51^{\mathrm m}\,35^{\mathrm s}$ & $216^\circ\,41'\,50''1$ & $0^\circ\,58'\,16''$ \\
\phantom{April }24 & $10\;\;32\;\;55$ & $216\;\;17\;\;58.7$ & $0\;\;46\;\;50$ \\
\phantom{April }25 & $10\;\;19\;\;22$ & $216\;\;6\;\;3.6$ & $0\;\;40\;\;4$ \\
\phantom{April }26 & $10\;\;17\;\;30$ & $215\;\;54\;\;7.3$ & $0\;\;33\;\;51$ \\
\end{tabular}
\end{center}

\SourceLeaf{336}\Needspace{12\baselineskip}\par\indent Die Declination vom 22. April ist zweifelhaft, und die vom 24. auch nicht so zuverlässig als die beiden folgende, wo der Planet eine bequemere Lage hatte. Für die Opposition hat Herr Prof. \textsc{Gauss} aus diesen Beobachtungen folgendes Resultat gefunden:
\begin{center}\fontsize{10}{12}\selectfont\setlength{\tabcolsep}{4pt}\renewcommand{\arraystretch}{1.15}
\begin{tabular}{l}
$\text{\AstroSymbols ☍}$ 1811. 24. April $19^{\mathrm u}\,20^{\mathrm m}\,12^{\mathrm s}$ Mittlere Zeit in Göttingen \\
$214^\circ\,8'\,48''3$ wahre heliocentrische Länge \\
$12\;\;55\;\;2.0$ wahre geocentrische Breite nordl. \\
\end{tabular}
\end{center}


\Needspace{20\baselineskip}\par\indent In der starken Abweichung dieses Ortes von den auf die vier ersten Oppositionen gegründeten Elementen, erkennt man nunmehr auch bei der Juno den Einfluss der Störungen, welche besonders der \emph{Jupiter} ausübt und deren Berechnung bei der Juno eine eben so ungeheure Arbeit erfordern wird, wie bei der Pallas; es ist billig, dass diese Arbeit bei der früher entdeckten Pallas zuerst beseitigt sein müsse, daher Herr Prof. \textsc{Gauss} sich einstweilen begnügt hat, neue elliptische Elemente auf die Oppositionen von 1806, 1807, 1808, 1811 zu gründen, welche wir hier folgen lassen:

\begin{center}\fontsize{10}{12}\selectfont\setlength{\tabcolsep}{4pt}\renewcommand{\arraystretch}{1.15}\begin{tabular}{llr}
\shortstack[l]{Epoche der mittlern Länge 1811\\Meridian von Göttingen} & .\;.\;.\;.\;. & $177^\circ\,48'\,1''8$\\
tägliche mittlere tropische Bewegung &  & $813''2486$\\
Länge der Sonnennähe 1811 & .\;.\;. & $53^\circ\,14'\,32''4$\\
Länge des aufsteigenden Knoten 1811 &  & $171^\circ\,9'\,13''5$\\
Neigung der Bahn & .\;.\;.\;.\;.\;.\;. & $13\;4\;27.0$\\
Excentricität & .\;.\;.\;.\;.\;. & $=\mathrm{Sin}\;14\;44\;9.1$\\
Log. der halben grossen Axe &  & $=0.4265711.$\\
\end{tabular}\end{center}

\par\indent Nach diesen Elementen hat Herr \textsc{Wachter}, ein geschickter Schüler des Herrn Professor \textsc{Gauss}, folgende Ephemeride für den Lauf der \emph{Juno} im Jahre 1812 berechnet. [vom 23. Febr. bis 5. Novemb.]
\clearpage\fancyhead[L]{\small PALLAS.}\hypertarget{native-336-lower}{}\label{native-336-lower}
\begin{center}\textbf{PALLAS.}\par\vspace{4pt}\hrule\vspace{4pt}\fontsize{10}{12}\selectfont Astronomisches Jahrbuch für 1814. Seite 246\,..\,249. Berlin 1811.\par\vspace{4pt}\hrule\vspace{4pt}\end{center}

\begin{flushright}Göttingen den 10. Februar 1811.\end{flushright}
\par\indent Unterm 14. Jan. c. habe ich in unsere gel. Anzeigen eine kurze Nachricht über die Störungen der Pallas vom Jupiter einrücken lassen. Es ist dies freilich noch nicht die \emph{allgemeine} Theorie jener Störungen; diese wird erst noch viel mehr Arbeit kosten, wenn sie alles, was zu wünschen ist, leisten soll. Ich habe damit auch schon einen Anfang gemacht, indessen pressirt es damit noch nicht, da die Art, wie ich die Störungen in meiner neulichen Arbeit behandelt habe, sich noch viele Jahre eben so fortsetzen lässt. Sobald die Opposition des gegenwärtigen Jahres beobachtet ist, werde ich diese Fortsetzung unternehmen, vielleicht dann auf eine ähnliche Art auch die Juno und Vesta behandeln. Ich ersuche Sie, die $\text{\AstroSymbols ☍}$ der Pallas, die nächstens eintritt, sorgfältig zu beobachten; zur Erleichterung hier die verbesserte Ephemeride, denn die Harmonie meiner neuen Resultate hat sich vollkommen bestätigt.
\Needspace{14\baselineskip}
\begin{center}\fontsize{10}{12}\selectfont\setlength{\tabcolsep}{4pt}\renewcommand{\arraystretch}{1.15}\begin{tabular}{lll}
Febr. 11 & AR. $148^\circ\,3'$ & Decl. $15^\circ\,43'$ S.\\
\phantom{Febr. }15 & $147\;\;20$ & $14\;\;16$\\
\phantom{Febr. }19 & $146\;\;37$ & $12\;\;42$\\
\phantom{Febr. }23 & $145\;\;55$ & $11\;\;3$\\
\phantom{Febr. }27 & $145\;\;17$ & $9\;\;19$\\
März \phantom{1}3 & AR. $144^\circ\,43'$ & Decl. $7^\circ\,32'$ S.\\
\phantom{März }7 & $144\;\;13$ & $5\;\;44$\\
11 & $143\;\;49$ & $3\;\;57$\\
15 & $143\;\;31$ & $2\;\;11$\\
19 & $143\;\;20$ & $0\;\;29$\\
\end{tabular}\end{center}

\SourceLeaf{337}\Needspace{12\baselineskip}\par\indent Prof. \textsc{Harding} hat in Manheim folgende Beobachtungen der Pallas gemacht:
\begin{center}\fontsize{10}{12}\selectfont\setlength{\tabcolsep}{4pt}\renewcommand{\arraystretch}{1.15}
\begin{tabular}{l|c|c|c}
1811 Januar 20 & $14^{\mathrm u}\,6^{\mathrm m}\,38^{\mathrm s}6$ & $151^\circ\,10'\,36''2$ & $-21^\circ\,7'\,39''7$ \\
\phantom{1811 Januar }21 & $14\;\;2\;\;18.1$ & $151\;\;4\;\;53.5$ & $-20\;\;57\;\;27.7$ \\
\phantom{1811 Januar }22 & $13\;\;57\;\;58.4$ & & $-20\;\;47\;\;51.7$ \\
\phantom{1811 Januar }23 & $13\;\;53\;\;38.6$ & $150\;\;52\;\;48.2$ & $-20\;\;37\;\;58.8$ \\
\phantom{1811 Januar }24 & $13\;\;49\;\;16.4$ & $150\;\;46\;\;0.6$ & $-20\;\;27\;\;9.7$ \\
\phantom{1811 Januar }25 & $13\;\;44\;\;53.0$ & $150\;\;39\;\;9.9$ & $-20\;\;16\;\;9.4$ \\
\end{tabular}
\end{center}


\par\indent Der Planet ist jetzt ziemlich hell, doch schien mir die Ceres, die sehr nahe auf dem Platze ihrer Ephemeride steht, noch etwas mehr Licht zu haben.
\Needspace{18\baselineskip}\begin{flushright}Göttingen den 29. August 1811.\end{flushright}
\par\indent Beigehend, Verehrtester Freund, habe ich das Vergnügen, Ihnen die Ephemeriden für \emph{Pallas}, \emph{Juno} und \emph{Vesta} für ihre nächste Erscheinung zu übersenden. Drei geschickte Schüler von mir haben sich in diese Rechnung getheilt. In Rücksicht der Pallas finden Sie noch das Nähere in Nr. 8 und 130 der Göttingischen Gelehrten Anzeigen [S. 327 d. B] angeführt; Hr. \textsc{Nicolai} hat unter meiner Aufsicht die Störung der Pallas durch Jupiter jetzt noch ein Jahr weiter fortgesetzt, und ist jetzt beschäftigt, die Ephemeride für 1813 zu berechnen. Das Nähere in Betreff der Juno steht in Nr. 92 unserer gel. Anz. [Seite 325 d. B.] Die Elemente der Vesta sind erst vorläufig und werden noch einer kleinen Verbesserung bedürfen. Inzwischen setze ich Ihnen die hiesigen am Mauerquadranten von mir gemachten Beobachtungen her.
\Needspace{9\baselineskip}\begin{center}\fontsize{10}{12}\selectfont\setlength{\tabcolsep}{4pt}\renewcommand{\arraystretch}{1.15}
\begin{tabular}{l|c|c|c}
1811 Mai 20 & $12^{\mathrm u}\,27^{\mathrm m}\,45''1$ M.\,Z. & $244^\circ\,42'\,0''9$ & $-12^\circ\,28'\,10''1$ \\
\phantom{1811 Mai }24 & $12\;\;8\;\;3.4$ & $243\;\;42\;\;26.1$ & $-12\;\;30\;\;1.8$ \\
\phantom{1811 Mai }25 & $12\;\;3\;\;7.6$ & $243\;\;27\;\;25.3$ & $-12\;\;30\;\;47.8$ \\
\phantom{1811 Mai }29 & $11\;\;43\;\;22.9$ & $242\;\;27\;\;0.7$ & $-12\;\;35$ \\
\phantom{1811 Mai }30 & $11\;\;38\;\;27.3$ & $242\;\;12\;\;3.0$ & $-12\;\;35\;\;52.8$ \\
\end{tabular}
\end{center}


\par\indent Ich habe aus diesen Beobachtungen, welche ich alle für gut halte, die Opposition folgendermassen abgeleitet: 1811 Mai 15. $12^{\mathrm u}\,44^{\mathrm m}\,30^{\mathrm s}$ M.\,Z. in Göttingen, $243^\circ\,48'\,43''9$ wahre Länge, $8^\circ\,33'\,59''5$ geocentr. nordl. Breite. Was die Veränderungen betrifft, welche Hr. Dr. \textsc{Triesnecker} an meinen Resultaten einiger Ceresoppositionen nach dem Jahrbuche von 1813 anbringen zu müssen geglaubt hat, so bemerke ich

\par\indent 1) Dass ich in der Berechnung der Opposition für 1807 aus den Mailänder Beobachtungen, so weit ich solche bisher nachgesehen, keine Fehler gefunden habe, sondern bei meinem Resultate 3. Mai $8^{\mathrm u}\,42^{\mathrm m}\,28^{\mathrm s}$ Par Z. und $222^\circ\,14'\,9''1$ Länge bleiben muss; inzwischen will ich die ganze Rechnung nochmals wiederholen, sobald ich dazu kommen kann, die Elemente nach den beiden Oppositionen von 1809 und 1811 (die noch nicht benutzt sind), zu berechnen.

\par\indent 2) Die Conjectur des Hrn. Dr. \textsc{Triesnecker}, bei der Opposition von 1808 die geoc. Breite $12^\circ\,43'\,44''8$ in $12^\circ\,45'\,44''8$ zu verwandeln, ist durchaus unstatthaft und eben so muss die Länge der Knotenlinie nothwendig $80^\circ\,53'\,23''$ bleiben, und keinesweges um $5'$ grösser angenommen werden. Hr. Tr. ist zu seiner Conjectur lediglich dadurch veranlasst worden, dass die Elemente mit meinen Perturbationstafeln sonst nicht so gut mit der Breite stimmen, aber ich wundere mich sehr, dass er hierbei übersehen hat, was ich doch ausdrücklich bemerkt habe, dass bei Berechnung der Breiten nicht meine alten Tafeln, sondern die späterhin von mir entwickelten vollständigen Störungsgleichungen angewandt sind, welche bisher noch nicht gedruckt wurden, und die ich deshalb Ihnen, so wie sie im Jahre 1805 berechnet sind, hier mittheile.
\SourceLeaf{338}\Needspace{24\baselineskip}\begin{center}\emph{Störungen der Breite der Ceres durch Jupiter.}

($\JupiterSymbol$ mittl. Länge des $\JupiterSymbol$, $\CeresSymbol$ mittl. Länge der $\CeresSymbol$)
\end{center}
\[
\begin{array}{l}
-9''65\,\sin(\JupiterSymbol-74^\circ\,13')\\
-14.17\,\sin(\CeresSymbol-2\JupiterSymbol+78^\circ\,35')\\
+27.47\,\sin(2\CeresSymbol-3\JupiterSymbol+78\;35)\\
-4.62\,\sin(3\CeresSymbol-4\JupiterSymbol+78\;35)\\
+5.51\,\sin(2\CeresSymbol-\JupiterSymbol-78\;35)\\
+0.99\,\sin(3\CeresSymbol-2\JupiterSymbol-78\;35)\\
+1.91
\end{array}
\]
\[
\begin{array}{l}
+8.61\,\sin(\CeresSymbol-\JupiterSymbol+95\;26)\\
-10.31\,\sin(2\CeresSymbol-2\JupiterSymbol+94\;10)\\
-6''16\,\sin(3\CeresSymbol-3\JupiterSymbol+112^\circ\,26')\\
-13.69\,\sin(2\JupiterSymbol-57^\circ\,21')\\
-6.77\,\sin(3\JupiterSymbol-\CeresSymbol-61\;25)\\
+4.30\,\sin(2\CeresSymbol-4\JupiterSymbol+72\;24)\\
-22.51\,\sin(3\CeresSymbol-5\JupiterSymbol+70\;27)\\
+1.65\,\sin(3\JupiterSymbol-56^\circ\,9')
\end{array}
\]
\par\indent Die sechs ersten Gleichungen sind die alten, aber verbessert, hingegen die neun übrigen, die Eine Minute betragen können, sind ganz neu hinzugekommen, und es ist daher nicht zu verwundern, dass bei Zuziehung aller die Lage des Knotens eine Veränderung von $5'$ erlitt. — Diese 15 Breitengleichungen sind noch zu ertragen, aber bei der Pallas habe ich über 40, die über $1''$ gehen.\par
\indent Den Cometen haben wir am 22. Aug. zuerst wiedergesehn, allein erst heute früh hat Prof. \textsc{Harding} einige Abstände desselben von Fixsternen gemessen, welche aber noch nicht reducirt sind.
